\documentclass[11pt,a4paper,UTF8]{ctexart}
\usepackage[margin=22mm,headheight=14pt]{geometry}
\usepackage{amsmath,amssymb,mathtools,bm,booktabs,longtable,array}
\usepackage{xcolor,listings}
\usepackage[most]{tcolorbox}
\usepackage{fancyhdr,hyperref,enumitem,microtype}

\definecolor{navy}{HTML}{12304A}
\definecolor{teal}{HTML}{176B87}
\definecolor{lightblue}{HTML}{EEF8FB}
\definecolor{codebg}{HTML}{F3F5F7}
\hypersetup{colorlinks=true,linkcolor=teal,urlcolor=teal}
\setlength{\parskip}{0.55em}
\setlength{\parindent}{2em}
\setlist[itemize]{leftmargin=2em,itemsep=0.25em}
\setlist[enumerate]{leftmargin=2.2em,itemsep=0.3em}
\renewcommand{\arraystretch}{1.25}
\newcolumntype{L}[1]{>{\raggedright\arraybackslash}p{#1}}
\newcommand{\pkg}[1]{\texttt{#1}}
\lstset{basicstyle=\ttfamily\small,backgroundcolor=\color{codebg},frame=single,
rulecolor=\color{teal!40},breaklines=true,columns=fullflexible,
xleftmargin=0.5em,xrightmargin=0.5em,aboveskip=0.7em,belowskip=0.7em}
\pagestyle{fancy}
\fancyhf{}
\lhead{\small OptoSimLab | 可微分光信号级仿真平台}
\rhead{\small 使用说明与自检记录}
\cfoot{\small 第 \thepage\ 页}
\renewcommand{\headrulewidth}{0.4pt}

\title{\textcolor{navy}{\Huge OptoSimLab}\\[0.6em]
\textcolor{teal}{\Large 使用说明与公式自检记录}}
\author{版本 0.8.0 \quad 项目位置：\texttt{F:\textbackslash OptoSimLab}}
\date{2026-09-06}

\begin{document}
\maketitle
\begin{tcolorbox}[colback=lightblue,colframe=teal,title=本次重制说明]
本手册使用 \LaTeX{} 原生数学环境排版全部公式。中文正文与数学公式分离渲染，不再把希腊字母、上下标和公式符号作为普通文本输出；本文同时是使用说明和当前版本的自检记录。
\end{tcolorbox}
\tableofcontents
\clearpage

\section{平台范围、交付与边界}
OptoSimLab 是基于 PyTorch 的\textbf{复包络、信号级、可微分}光学仿真平台。它用于搭建光通信链路、可训练光计算模块和小型光神经网络；它不求解波导横向模式、纳米尺度电磁场或 FDTD 场分布。

全部器件按“一器件一 \pkg{torch.nn.Module}”封装。需要优化的物理量注册为 \pkg{nn.Parameter}，可被优化器、\pkg{state\_dict()} 和 \pkg{.to(device)} 管理。

\begin{longtable}{L{27mm}L{82mm}L{44mm}}
\toprule
\textbf{类别} & \textbf{已实现内容} & \textbf{验证重点}\\
\midrule
\endhead
核心约定 & \pkg{SimulationGrid}、FFT 频率轴、物理常数 & 频率轴、单位\\
调制器 & 理想/有限消光比 MZM、低通驱动 MZM、实测曲线 MZM & 峰值、插值、梯度\\
频域器件 & Butterworth 低通、WaveShaper、色散光纤 & 截止点、损耗、相位\\
光梳 & 严格 FFT bin 对齐的频域光梳 & 谱线位置、总功率\\
电光梳 & 理想级联相位调制、RF 一阶响应、Wiener 相位噪声 & Bessel 边带、功率\\
放大器 & 商用增益控制 EDFA；旧电流模型仅兼容 & 带宽、饱和、ASE\\
干涉器件 & 1$\times$2 分束器、2$\times$2 耦合器、单/双输入 MZI & 幺正性、差分响应\\
功能模块 & 串行/图式器件链、Nature 2021 实值光卷积 & 分支汇合、逐点卷积\\
系统指标 & 频谱、眼图、EVM、Gaussian-Q BER & Parseval、一致性\\
\bottomrule
\end{longtable}

\begin{tcolorbox}[colback=lightblue,colframe=teal,title=真实器件边界]
当前是参数化工程模型。MZM 已支持实测 CSV 曲线导入和插值，但仓库不附带特定硬件的测量数据。没有可追溯的实测数据时，不能把仿真结果宣称为已与任一真实设备完全一致。
\end{tcolorbox}

\section{环境、安装和首次验证}
\subsection{进入已有环境}
在 PowerShell 中执行。若提示符开头出现 \pkg{(.venv)}，代表虚拟环境已经激活。
\begin{lstlisting}[language=bash]
cd F:\OptoSimLab
.\.venv\Scripts\Activate.ps1
python -m pytest -q -p no:cacheprovider tests
python examples\train_mzm_bias.py
\end{lstlisting}
当前 Windows CPU 环境的验收结果是：\textbf{82 项通过，1 项跳过}。跳过项是 CPU/CUDA 一致性测试，因为当前 PyTorch 不带 CUDA；这不是测试失败。

\subsection{从零安装}
\begin{lstlisting}[language=bash]
cd F:\OptoSimLab
python -m venv .venv
.\.venv\Scripts\Activate.ps1
python -m pip install --upgrade pip
python -m pip install torch numpy pytest
python -m pip install -e . --no-deps
python -m pytest -q -p no:cacheprovider tests
\end{lstlisting}

\subsection{安装已打包版本}
\begin{lstlisting}[language=bash]
cd F:\OptoSimLab
.\.venv\Scripts\Activate.ps1
python -m pip install --force-reinstall --no-deps .\dist\optosimlab-0.8.0-py3-none-any.whl
python -c "from optosimlab import PhotonicRealConvolution; print('import OK')"
\end{lstlisting}

\section{统一信号与数学约定}
\subsection{复包络、采样和 FFT}
时域光场采用复包络 $E(t)$，频率 $f$ 是相对中心载波的偏移，不直接对数百 THz 的载波采样。PyTorch DFT 约定为
\begin{align}
X[k] &= \sum_{n=0}^{N-1}x[n]\exp\!\left(-\mathrm{j}\frac{2\pi kn}{N}\right),\\
x[n] &= \frac{1}{N}\sum_{k=0}^{N-1}X[k]\exp\!\left(+\mathrm{j}\frac{2\pi kn}{N}\right),\\
\Delta f &= \frac{f_{\mathrm{s}}}{N}.
\end{align}
器件必须从 \pkg{SimulationGrid} 获得频率轴和 $\Delta f$，而不应各自复制采样率常数。

\subsection{功率、损耗和频谱}
\begin{align}
P[n]&=|E[n]|^2, &
\bar P&=\frac{1}{N}\sum_{n=0}^{N-1}|E[n]|^2,\\
\eta_P&=10^{-L_{\mathrm{dB}}/10}, &
\eta_E&=\sqrt{\eta_P}=10^{-L_{\mathrm{dB}}/20}.
\end{align}
这里 $L_{\mathrm{dB}}$ 是功率损耗。将 $10^{-L_{\mathrm{dB}}/10}$ 直接乘到场上会使功率损耗加倍，是本项目的重点公式自检项。频谱采用
\begin{equation}
P_{\mathrm{bin}}[k]=\frac{|X[k]|^2}{N^2},\qquad
\sum_kP_{\mathrm{bin}}[k]=\bar P,
\end{equation}
所以满足 Parseval 一致性。

\subsection{默认物理单位}
公开接口统一采用 SI：光场为 $\sqrt{\mathrm{W}}$，光功率为 $\mathrm{W}$，电压为 $\mathrm{V}$，时间为 $\mathrm{s}$，频率为 $\mathrm{Hz}$，波长与长度为 $\mathrm{m}$，相位为 $\mathrm{rad}$。显示功率可使用 dBm：
\begin{align}
P_{\mathrm{W}}&=10^{P_{\mathrm{dBm}}/10}\times10^{-3},\\
P_{\mathrm{dBm}}&=10\log_{10}\!\left(\frac{P_{\mathrm{W}}}{10^{-3}}\right).
\end{align}
例如 $0\,\mathrm{dBm}=1\,\mathrm{mW}=10^{-3}\,\mathrm{W}$。源码中的 \pkg{units.py} 提供统一换算，禁止混用 mW 与 W。

\section{器件公式与正确使用方式}
\subsection{MZM 与电驱动低通}
理想推挽 MZM 的场传输公式为
\begin{equation}
E_{\mathrm{out}}(t)=E_{\mathrm{in}}(t)
\cos\!\left[\frac{\pi\bigl(V(t)+V_{\mathrm{bias}}\bigr)}{2V_{\pi}}\right].
\end{equation}
插损在场上使用 $10^{-L_{\mathrm{dB}}/20}$。有限消光比用正交泄漏项实现，峰谷\textbf{功率比}为
\begin{equation}
\frac{P_{\max}}{P_{\min}}=10^{\mathrm{ER}_{\mathrm{dB}}/10}.
\end{equation}
电驱动低通为零相位 Butterworth 幅度响应：
\begin{equation}
H_{\mathrm{LP}}(f)=\left[1+\left(\frac{|f|}{f_c}\right)^{2n}\right]^{-1/2}.
\end{equation}
因此在 $|f|=f_c$ 处有 $|H_{\mathrm{LP}}|^2=1/2$，即 $-3.0103\,\mathrm{dB}$ 功率点。

实测 MZM 用电压节点 $(V_i,T_i,\phi_i)$ 描述，$T_i\in[0,1]$ 是功率透射率。在相邻节点之间采用分段线性插值，超出标定范围时钳位到端点：
\begin{align}
T(V)&=(1-a)T_i+aT_{i+1},\\
\phi(V)&=(1-a)\phi_i+a\phi_{i+1},\qquad
a=\frac{V-V_i}{V_{i+1}-V_i},\\
E_{\mathrm{out}}&=E_{\mathrm{in}}\sqrt{T(V)}\exp[\mathrm{j}\phi(V)].
\end{align}
CSV 默认列为 \pkg{voltage\_v,power\_transmission}，可选列为 \pkg{phase\_rad}。钳位是为了避免对未测区域作无依据外推，不代表真实器件在范围外恒定。

\subsection{频域光梳与理想电光梳}
频域光梳每条谱线都必须落在 FFT bin 上：
\begin{align}
f_m &= m\Delta f,\\
E(t)&=\sum_m\sqrt{P_m}\exp\!\left[\mathrm{j}(2\pi f_mt+\phi_m)\right],\\
\bar P&=\sum_mP_m.
\end{align}
若线间隔不是 $\Delta f$ 的整数倍，有限记录会产生谱泄漏，实现会拒绝该输入。

多级相位调制的理想电光梳为
\begin{equation}
E_{\mathrm{out}}(t)=E_{\mathrm{in}}(t)
\exp\!\left[\mathrm{j}\sum_q\beta_q\sin(2\pi f_qt+\phi_q)\right]
10^{-L_{\mathrm{dB}}/20}.
\end{equation}
无插损时指数因子的模恒为 1，所以逐点功率守恒。单音连续波的第 $n$ 条边带满足
\begin{equation}
E_n\propto J_n(\beta)\exp(\mathrm{j}n\phi),\qquad f_n=nf_{\mathrm{RF}}.
\end{equation}

\subsection{实用 RF 电光梳}
每一级的 RF 一阶响应、等效调制深度和相位为
\begin{align}
H_{\mathrm{RF}}(f)&=\frac{1}{1+\mathrm{j}f/f_c},\\
\beta_{\mathrm{eff}}&=\beta|H_{\mathrm{RF}}|,\qquad
\phi_{\mathrm{eff}}=\phi+\arg H_{\mathrm{RF}}.
\end{align}
可选 Wiener 相位扩散满足
\begin{equation}
\operatorname{Var}\!\left[\varphi_{n+1}-\varphi_n\right]=D_{\varphi}\Delta t.
\end{equation}
这是可控的相位扩散近似，并非完整的真实激光器相噪模型。

\subsection{光纤与 WaveShaper}
线性色散光纤的频域场传递函数和参数换算为
\begin{align}
H_{\mathrm{fiber}}(f)&=
\exp\!\left(-\frac{\alpha_PL}{2}\right)
\exp\!\left(\mathrm{j}\frac{\beta_2(2\pi f)^2L}{2}\right),\\
\beta_2&=-\frac{D\lambda_0^2}{2\pi c}.
\end{align}
WaveShaper 在 \pkg{fftshift} 频率轴上插值可训练的衰减与相位控制点：
\begin{equation}
H_{\mathrm{WS}}(f)=10^{-A(f)/20}\exp\!\bigl(\mathrm{j}\theta(f)\bigr).
\end{equation}

\subsection{分束器、定向耦合器与 MZI}
1$\times$2 分束器的功率分配比为 $r$，无插损时
\begin{equation}
\begin{bmatrix}E_0\\E_1\end{bmatrix}
=\begin{bmatrix}\sqrt r\\\sqrt{1-r}\end{bmatrix}E_{\mathrm{in}},
\qquad |E_0|^2+|E_1|^2=|E_{\mathrm{in}}|^2.
\end{equation}
2$\times$2 定向耦合器保留交叉耦合的正交相位：
\begin{equation}
\bm C(r)=
\begin{bmatrix}
\sqrt r & \mathrm{j}\sqrt{1-r}\\
\mathrm{j}\sqrt{1-r} & \sqrt r
\end{bmatrix},\qquad \bm C^{\mathrm H}\bm C=\bm I.
\end{equation}
MZI 使用 $\bm C\,\operatorname{diag}(e^{\mathrm{j}x},e^{-\mathrm{j}x})\bm C$，其中
$x=\pi(V+V_{\mathrm{bias}})/(2V_\pi)$。当 $r=1/2$ 且观察端口为 1 时，双输入单输出结果为
\begin{equation}
E_{\mathrm{out}}=\mathrm{j}\left[\cos(x)E_0-\sin(x)E_1\right].
\end{equation}
令 $E_1=0$ 即得到单输入单输出 MZI。旧 \pkg{FourInputMZI} 仅为 v0.7 代码兼容，不再作为推荐物理模型。

\subsection{EDFA、饱和与 ASE}
推荐的 \pkg{GainControlledEDFA} 使用商用放大器常见的净增益设定 $G_{\mathrm{set,dB}}$，而不是让用户控制内部泵浦电流。未饱和功率增益为
\begin{align}
G_{\mathrm{ss}}(f)&=10^{G_{\mathrm{set,dB}}/10}
\exp\!\left[-4\ln2\left(\frac{f-f_0}{B}\right)^2\right].
\end{align}
在 $f=f_0\pm B/2$ 处功率增益为峰值的一半；场使用 $\sqrt{G_{\mathrm{ss}}(f)}$。

以输出饱和功率 $P_{\mathrm{sat,out}}$ 表示的平滑压缩模型为
\begin{align}
G_{\mathrm{eff}}&=1+\frac{G_{\mathrm{ss}}-1}
{1+G_{\mathrm{ss}}P_{\mathrm{in}}/P_{\mathrm{sat,out}}},
\qquad G_{\mathrm{ss}}\ge 1.
\end{align}
当 $G_{\mathrm{ss}}<1$ 时保持原频带整形值，不把带外损耗错误改为单位增益。一偏振复包络、双边 PSD 下的 ASE 为
\begin{align}
S_{\mathrm{ASE}}(f)&=n_{\mathrm{sp}}h\nu
\max(G_{\mathrm{eff}}(f)-1,0),&
n_{\mathrm{sp}}&=\frac{10^{\mathrm{NF}_{\mathrm{dB}}/10}}{2}.
\end{align}
对长度为 $N$ 的 \pkg{ifft}，频域复高斯噪声方差取
\begin{equation}
\operatorname{Var}\{N[k]\}=N^2S_{\mathrm{ASE}}(f_k)\Delta f,
\end{equation}
于是 IFFT 后的期望平均噪声功率为 $\sum_kS_{\mathrm{ASE}}(f_k)\Delta f$。

旧 \pkg{SmallSignalEDFA} 和 \pkg{SaturatedNoisyEDFA} 的泵浦电流接口仅为 v0.7 兼容，不应用于新的商用品级链路配置。

\section{功能模块、链路和训练}
\subsection{复值卷积和长序列 FIR}
\pkg{OpticalChain} 保留串行构造方式，并新增 \pkg{from\_graph}。图节点显式声明模块、输入名和输出名，按拓扑顺序执行，因此同一个值可扇出到多个并行支路，分束器可返回多个端口，耦合器或探测器可重新汇合多个输入。若只有一个最终输出则返回张量；多个最终输出按名称返回字典。

\pkg{ComplexConv1d} 的实虚部由四次实卷积组成：
\begin{align}
\Re(y)&=x_r*w_r-x_i*w_i,\\
\Im(y)&=x_r*w_i+x_i*w_r.
\end{align}
偏置只加入一次，避免重复。
它是数学复值网络兼容层，并不表示 Xu 等人的微梳光卷积硬件；物理路线应使用 \pkg{PhotonicRealConvolution}。

\pkg{OverlapSaveFIR} 实现可训练复因果 FIR：
\begin{equation}
y[n]=\sum_{m=0}^{M-1}h[m]x[n-m],\qquad
N_{\mathrm{FFT}}>M-1.
\end{equation}
它在输入前补零，将每个频域块开头 $M-1$ 个循环卷积样本丢弃，因此输出严格对应线性因果卷积。

\subsection{Nature 2021 实值 WDM 光卷积}
实现链路遵循 Xu 等人在 Nature 2021 报道的结构：等间隔光梳在一根光纤中并行传播，WaveShaper 将卷积核系数映射到各梳齿功率，单个强度调制器把待计算序列广播到所有梳齿，色散使第 $k$ 个通道产生 $kD$ 个采样点的递增延时，最后在同一时刻平方律探测并求和。

设每条输入梳齿功率均为 $P_{\mathrm{line}}$，单支路权重 $w_k\in[0,1]$，归一化输入 $x[n]\ge0$。WaveShaper 和广播调制后的通道场分别为
\begin{align}
E_k^{(w)}[n]&=\sqrt{w_k}\,E_k[n],\\
E_k^{(m)}[n]&=\sqrt{x[n]},E_k^{(w)}[n].
\end{align}
零填充递增延时与低带宽探测器输出为
\begin{align}
E_k^{(d)}[n]&=E_k^{(m)}[n-kD],\\
P_{\mathrm{det}}[n]&=\sum_k\left|E_k^{(d)}[n]\right|^2
=P_{\mathrm{line}}\sum_k w_kx[n-kD].
\end{align}
因此校准输出
\begin{equation}
y[n]=\frac{P_{\mathrm{det}}[n]}{P_{\mathrm{line}}}
=\sum_k w_kx[n-kD]
\end{equation}
逐点等于零初始状态的因果实卷积。探测器“先按通道求功率再相加”等效于其电带宽低于梳齿间隔，因而滤除通道间拍频。单个无源 WaveShaper 支路只能给出同号权重。

有符号卷积使用 \pkg{DifferentialPhotonicConvolution}，将正、负权重分到两个探测支路后作差。

\begin{lstlisting}[language=Python]
import torch
from optosimlab import PhotonicRealConvolution

x = torch.tensor([0.1, 0.3, 0.7, 1.0, 0.2, 0.6])
w = torch.tensor([0.2, 0.5, 0.8])
conv = PhotonicRealConvolution(
    w, sample_rate_hz=100e9, line_spacing_hz=10e9,
    line_power_w=1e-3, tap_delay_samples=1)
y = conv(x)
p_w = conv.detected_power_w(x)
assert torch.allclose(p_w, 1e-3 * y, atol=1e-9)
\end{lstlisting}

\subsection{最小可运行链路}
\begin{lstlisting}[language=Python]
import torch
from optosimlab import SimulationGrid, MachZehnderModulator
from optosimlab import LinearDispersiveFiber

grid = SimulationGrid(sample_rate_hz=100e9)
samples = 4096
t = torch.arange(samples) / grid.sample_rate_hz
carrier = torch.ones(samples, dtype=torch.complex64)
drive = 0.5 * torch.sin(2 * torch.pi * 5e9 * t)
mzm = MachZehnderModulator(v_pi=1.0, bias_voltage=0.5)
fiber = LinearDispersiveFiber(
    grid, length_m=10e3, attenuation_db_per_km=0.2,
    dispersion_ps_nm_km=17.0)
field_out = fiber(mzm(carrier, drive))
\end{lstlisting}

\subsection{只训练需要训练的参数}
\begin{lstlisting}[language=Python]
optimizer = torch.optim.Adam([mzm.bias_voltage], lr=3e-2)
optimizer.zero_grad()
field = mzm(carrier, drive)
loss = (field.abs().square().mean() - 0.5).square()
loss.backward()
optimizer.step()
\end{lstlisting}
可训练不代表所有参数都必须训练；只把有明确物理含义且允许调节的参数交给优化器。

\section{自检记录与验收门禁}
\subsection{已完成的专项自检}
\begin{longtable}{L{44mm}L{110mm}}
\toprule
\textbf{检查对象} & \textbf{已验证内容}\\
\midrule
\endhead
MZM & 峰值/零点；$3\,\mathrm{dB}$ 插损的场与功率关系；$20\,\mathrm{dB}$ 消光比；梯度。\\
实测 MZM & 标定节点、区间线性插值、范围外钳位、CSV 列读取及曲线参数梯度。\\
分束/耦合/MZI & 总功率守恒、2$\times$2 幺正性、交叉端口 $\mathrm{j}$ 相位、单/双输入解析峰值与零点。\\
低通与光纤 & Butterworth 半功率截止点；$D\leftrightarrow\beta_2$ 单位换算；损耗；无损色散能量守恒。\\
频域光梳与电光梳 & FFT bin、谱线总功率、Bessel 边带、RF 一阶响应、Wiener 相位扩散、非法输入。\\
EDFA & dB 增益设定、$-3\,\mathrm{dB}$ 带宽、饱和单调性、ASE PSD 非负和梯度；确认新接口无泵浦电流。\\
图式链路 & 分束器多输出、并行支路、双输入汇合、缺失连接错误提示。\\
Nature 实值卷积 & WaveShaper 功率权重、广播、零填充递增延时、探测功率及正/负双支路逐点对照直接卷积。\\
复值和长序列模块 & 复卷积偏置一次性计入；Overlap-Save 与直接线性卷积、批量维、梯度一致。\\
系统指标 & Parseval 频谱、眼图、EVM、等先验实高斯二元假设下的 Q 与 BER。\\
设备一致性 & CPU 路径通过；CUDA PyTorch 环境下自动执行 CPU/CUDA 一致性测试。\\
\bottomrule
\end{longtable}

\subsection{新增模块必须遵循的流程}
\begin{enumerate}
\item 写清输入/输出形状、实数或复数 dtype、单位及 FFT 定义。
\item 将可调物理量封装为 \pkg{nn.Parameter}，拒绝不合理输入。
\item 为至少一个可解析结论写测试，例如截止点、极值或功率守恒。
\item 为错误 shape、dtype、负参数及冲突输入写失败测试。
\item 确认目标参数的 \pkg{grad} 非空；高风险公式增加有限差分对照。
\item 先跑 CPU，再在有 CUDA 时比较 CPU/GPU；最后全量测试、打包并做导入测试。
\end{enumerate}

\subsection{复现实验命令}
\begin{lstlisting}[language=bash]
cd F:\OptoSimLab
.\.venv\Scripts\Activate.ps1
python -m compileall -q src tests examples
python -m pytest -q -p no:cacheprovider tests
python -m pip wheel --no-deps --wheel-dir dist .
python -c "from optosimlab import OverlapSaveFIR; print('package OK')"
\end{lstlisting}
\begin{tcolorbox}[colback=lightblue,colframe=teal,title=当前验收结果]
版本 0.8.0 已完成源码编译、专项测试、全量测试和 wheel 导入验证。全量结果：\textbf{82 passed, 1 skipped}。无 CUDA 的 CPU 环境中该跳过项是预期行为；在 CUDA 环境中它会自动运行。
\end{tcolorbox}

\section{当前完成度、限制与后续实验标定}
\begin{enumerate}
\item 软件层已经支持 MZM CSV 曲线；下一步应导入指定型号器件的真实电压、功率透射率与相位数据。
\item 应使用仪器数据标定 EDFA 增益谱、输出饱和功率和噪声系数；当前模型是可微的系统级近似，不含粒子数反转动力学。
\item 使用独立的保留测量曲线验证拟合误差，不能只报告训练曲线误差。
\item 当前递增延时由整数采样点精确实现；若实验色散对应分数采样延时，应新增带抗混叠设计的分数延时滤波器并单独验证。
\item 论文路线的通道功率求和假定探测器带宽低于梳齿间隔；若接收机能看到拍频，必须改用全场平方后再建模电接收带宽。
\end{enumerate}
\begin{tcolorbox}[colback=lightblue,colframe=teal,title=结论]
本平台已完成初始清单和 Issue \#1 所要求的软件功能、PyTorch 封装、解析公式验证、回归测试、打包和说明文档。这里的“完成”是指模型与接口完成，不等于已替代实验：下一阶段应使用可追溯的真实器件数据做标定和系统级一致性验证。
\end{tcolorbox}

\section{参考文献}
X.~Xu, M.~Tan, B.~Corcoran, et al., ``11 TOPS photonic convolutional accelerator for optical neural networks,'' \textit{Nature}, vol.~589, pp.~44--51, 2021. DOI: \href{https://doi.org/10.1038/s41586-020-03063-0}{10.1038/s41586-020-03063-0}.
\end{document}
